\documentclass[11pt,letterpaper]{article}

% ── Packages ──────────────────────────────────────────────────────────────────
\usepackage[margin=1in]{geometry}
\usepackage{amsmath,amssymb}
\usepackage{booktabs}
\usepackage{microtype}
\usepackage{parskip}
\usepackage{xcolor}
\usepackage{listings}
\usepackage{enumitem}
\usepackage{caption}
\usepackage{float}
\usepackage[colorlinks=true,linkcolor=blue,citecolor=blue,urlcolor=blue]{hyperref}

% ── Metadata ──────────────────────────────────────────────────────────────────
\title{%
  \textbf{AGI Pragma: A Decision Intelligence Core for Safe Agentic AI}%
}

\author{%
  Rafal Zabinski\\[4pt]
  Independent Research\\
  \texttt{github.com/zabinskirafal/AGI-Pragma-Core}
}

\date{}

% ── Code style ────────────────────────────────────────────────────────────────
\lstset{
  basicstyle=\ttfamily\small,
  frame=single,
  backgroundcolor=\color{gray!8},
  breaklines=true,
  captionpos=b,
}

% ══════════════════════════════════════════════════════════════════════════════
\begin{document}
\maketitle

% ── Abstract ──────────────────────────────────────────────────────────────────
\begin{abstract}
We present AGI Pragma, a seven-stage Decision Intelligence Core (DIC) designed
to enforce safety constraints and structured risk evaluation before agentic AI
systems execute actions. The DIC operates as a governance layer between any
action generator and the environment: it filters candidate actions through
branching, Monte Carlo critical path estimation, FMEA risk scoring, a
configurable decision gate, a four-state circuit breaker, utility-based
selection, and Bayesian belief update. We evaluate the system on three benchmark
environments --- Snake (10$\times$10 grid, self-collision), Maze (15$\times$15
recursive backtracker, dead-end traps), and Dynamic Threat Gridworld
(15$\times$15, five wandering hazards) --- comparing against a random baseline
that shares the same action filter. The DIC achieves $+$780\% score improvement
in Snake, 100\% solve rate in Maze (vs.\ 6\% baseline), and 78\% solve rate in
Gridworld (vs.\ 0\% baseline). We further implement and evaluate persistent
episodic memory via Bayesian prior transfer across sessions, document the
conditions under which it provides measurable lift, and derive four empirical
lessons about DIC calibration in practice.
\end{abstract}

% ── 1. Introduction ───────────────────────────────────────────────────────────
\section{Introduction}
\label{sec:intro}

Agentic AI systems --- language models with tool access, autonomous planners,
multi-agent pipelines --- are increasingly capable of consequential real-world
action: executing code, modifying files, calling external APIs, interacting with
live infrastructure. Their failure modes differ qualitatively from those of
passive models. A hallucinated output is correctable; a dropped database table
or a committed security misconfiguration may not be.

Existing safety approaches operate at the wrong level of abstraction. Post-hoc
monitoring detects failures after execution. Constitutional AI and RLHF shape
preferences but provide no enforcement. Prompt engineering nudges behavior without
hard constraints. None of these mechanisms answers the question that matters in a
deployed agentic context: \textit{is this specific action safe to execute right
now, given the current state of the environment?}

The gap is not the absence of risk knowledge --- it is the absence of a structured
mechanism that converts risk knowledge into a binding pre-execution constraint.

We address this gap with the Decision Intelligence Core (DIC), a seven-stage
pipeline that sits between action generation and action execution. The DIC makes
no assumptions about the upstream generator: an LLM, a rule-based planner, or a
learned policy may sit above it. Its role is not to generate better actions but
to govern which generated actions are permitted to execute, and to produce an
auditable record of every decision.

This paper makes the following contributions:
\begin{enumerate}[leftmargin=*, nosep]
  \item A fully specified seven-stage DIC architecture with domain-agnostic
        safety properties and domain-specific utility extension points.
  \item Empirical validation across three environments with increasing risk
        complexity, against a controlled random baseline sharing the same
        action filter.
  \item Implementation and experimental evaluation of cross-session Bayesian
        episodic memory, with a characterisation of the conditions under which
        it provides measurable lift.
  \item Four empirically-derived lessons about DIC calibration applicable to
        any agentic deployment context.
\end{enumerate}

% ── 2. Method ─────────────────────────────────────────────────────────────────
\section{Method: The Decision Intelligence Core}
\label{sec:method}

The DIC processes one decision per time step. Given the current environment state
and a set of candidate actions, it produces a single approved action and a complete
audit trace. Figure~\ref{fig:architecture} shows the system position; the seven
stages are described below.

\begin{figure}[H]
\centering
\begin{verbatim}
  ┌─────────────────────────────┐
  │   Action Generator          │  LLM, planner, neural policy
  └────────────┬────────────────┘
               │ candidate actions
               ▼
  ┌─────────────────────────────┐
  │   Decision Intelligence     │  Branching → Critical Path → FMEA
  │   Core  (DIC)               │  → Gate → Circuit Breaker
  │                             │  → Selection → Belief Update
  └────────────┬────────────────┘
               │ approved action + audit trace
               ▼
  ┌─────────────────────────────┐
  │   Environment / Execution   │  real world, simulation, API
  └─────────────────────────────┘
\end{verbatim}
\caption{DIC position in the system. The governance layer is indifferent to
the upstream generator.}
\label{fig:architecture}
\end{figure}

\paragraph{Stage 1 --- Branching.}
Candidate actions are enumerated and any action leading to an immediately fatal
or physically impossible state is removed before evaluation. This stage is shared
with the random baseline in all experiments; measured performance improvements are
attributable to Stages 2--7 only.

\paragraph{Stage 2 --- Critical Path Estimation.}
Each surviving candidate is evaluated by Monte Carlo rollout. From the state
that would result from the candidate action, $N$ trajectories of depth $d$ are
simulated under a lightweight random policy on a cloned environment. Each
rollout records whether the trajectory terminates in failure or depth exhaustion
(treated as timeout). Per-candidate estimates are produced:
\begin{itemize}[nosep]
  \item $p_{\mathrm{death}}$: fraction of rollouts ending in catastrophic failure,
  \item $p_{\mathrm{trap}}$: fraction where $\geq 3$ of 4 neighbours are
        simultaneously occupied (irreversible state),
  \item $\mathbb{E}[\text{steps to failure}]$: mean steps until first terminal event.
\end{itemize}
All experiments use $N = 200$ rollouts and depth $d = 25$ (Snake) or $d = 50$
(Maze, Gridworld).

\paragraph{Stage 3 --- FMEA Risk Scoring.}
Each candidate receives a Risk Priority Number:
\begin{equation}
  \mathrm{RPN} = \mathrm{Severity} \times \mathrm{Occurrence} \times \mathrm{Detection}
  \label{eq:rpn}
\end{equation}
Severity captures consequence magnitude (range 1--10). Occurrence is derived from
the Monte Carlo estimates. Detection captures how silently a failure mode manifests
--- silent failures receive Detection $= 1$ (highest penalty). Each domain defines
its own FMEA table. The maximum RPN across all failure modes is used for gating.

\paragraph{Stage 4 --- Decision Integrity Gate.}
Any candidate whose maximum RPN meets or exceeds a configurable threshold $\tau$
is blocked before utility evaluation. $\tau = 240$ in all experiments. The gate
is a hard constraint; utility estimates cannot override it.

\paragraph{Stage 5 --- Circuit Breaker.}
The circuit breaker maintains a four-state autonomy level based on the current
maximum RPN:

\begin{table}[H]
\centering
\begin{tabular}{lll}
\toprule
\textbf{State} & \textbf{RPN range} & \textbf{Effect} \\
\midrule
OK   & $< 180$   & No restriction \\
WARN & 180--219  & Decision depth logged \\
SLOW & 220--259  & Decision depth restricted \\
STOP & $\geq 260$ & Action blocked \\
\bottomrule
\end{tabular}
\caption{Circuit breaker states. When all candidates are blocked the system falls
back to the least-RPN candidate rather than failing open.}
\label{tab:circuit_breaker}
\end{table}

\paragraph{Stage 6 --- Decision Selection.}
Among unblocked candidates, the agent selects by utility. The general form is:
\begin{equation}
  U(a) = (1 - p_{\mathrm{death}}) \cdot w_1
       + (1 - p_{\mathrm{trap}}) \cdot w_2
       - \mathrm{dist} \cdot w_3
       - \mathrm{revisits} \cdot w_4
       - \frac{\mathrm{RPN}}{1000}
  \label{eq:utility}
\end{equation}
Domain-specific terms are added where appropriate: a proximity penalty in
Gridworld; exact BFS path distance in Maze.

\paragraph{Stage 7 --- Belief Update.}
After each action, Bayesian Beta trackers update hazard rate estimates from
observed risk signals. The posterior after episode $N$ is available as the
prior for episode $N+1$ via persistent episodic memory
(Section~\ref{sec:memory}).

% ── 3. Experiments ────────────────────────────────────────────────────────────
\section{Experiments}
\label{sec:experiments}

\subsection{Environments}

\textbf{Snake} is a 10$\times$10 grid where the agent navigates a growing snake
body toward food. Self-collision and wall collision are terminal. The dominant
risk is trap: the snake's own body creates irreversible configurations.

\textbf{Maze} is a 15$\times$15 grid generated by recursive backtracking. Wall
hits are no-ops. The agent navigates from the top-left cell to the bottom-right.
A BFS distance table is precomputed from the goal at episode start; the utility
function uses exact topological path distance, not Manhattan distance.

\textbf{Dynamic Threat Gridworld} is a 15$\times$15 grid with five wandering
hazard agents that move randomly each step. Contact with a hazard from any
direction is terminal. WAIT is a first-class action. This is the first
environment in the suite where $p_{\mathrm{death}}$ varies meaningfully across
candidate actions at every step --- moving toward a hazard cluster has
genuinely higher $p_{\mathrm{death}}$ than WAIT or evasion.

\subsection{Baseline}

The baseline agent selects uniformly at random from \texttt{safe\_actions()} ---
the same immediate-collision filter used by Stage 1 of the DIC. This isolates the
contribution of Stages 2--7. The baseline uses no FMEA, no Monte Carlo estimation,
no circuit breaker, and no utility function.

\subsection{Results}

All results are over 50 episodes per agent per benchmark, seeds 0--49.

\begin{table}[H]
\centering
\begin{tabular}{lrrr}
\toprule
\textbf{Metric} & \textbf{Baseline} & \textbf{DIC (Pragma)} & \textbf{$\Delta$} \\
\midrule
Avg score   & 2.6   & 22.5  & $+$780\% \\
Max score   & 7     & 33    & $+$371\% \\
Avg steps   & 284.2 & 198.0 & $-$30\%  \\
Avg reward  & 8.6   & 101.4 & $+$1081\% \\
\bottomrule
\end{tabular}
\caption{Snake --- 10$\times$10 grid, 50 episodes each. Lower steps is better
(fewer wasted steps before terminal state or food collection).}
\label{tab:snake}
\end{table}

\begin{table}[H]
\centering
\begin{tabular}{lrrr}
\toprule
\textbf{Metric} & \textbf{Baseline} & \textbf{DIC (Pragma)} & \textbf{$\Delta$} \\
\midrule
Solved              & 3/50 (6\%)   & 50/50 (100\%) & $+$1567\% \\
Avg steps (all)     & 296.1        & 46.1          & $-$84\%   \\
Avg steps (solved)  & 234.7        & 46.1          & $-$80\%   \\
Avg reward          & $-$38.3      & 5.5           & $+$114\%  \\
\bottomrule
\end{tabular}
\caption{Maze --- 15$\times$15 recursive backtracker, 50 episodes each.
The random baseline spent 234.7 steps on the 3 episodes it did solve;
Pragma solved all 50 in 46.1 steps on average.}
\label{tab:maze}
\end{table}

\begin{table}[H]
\centering
\begin{tabular}{lrrr}
\toprule
\textbf{Metric} & \textbf{Baseline} & \textbf{DIC (Pragma)} & \textbf{$\Delta$} \\
\midrule
Solved          & 0/50 (0\%)  & 39/50 (78\%) & ---       \\
Killed by hazard & 50/50      & 11/50        & $-$78\%   \\
Timed out       & 0/50        & 0/50         & ---       \\
Avg steps       & 56.8        & 22.8         & $-$60\%   \\
Avg reward      & $-$17.4     & 3.4          & $+$120\%  \\
\bottomrule
\end{tabular}
\caption{Dynamic Threat Gridworld --- 15$\times$15, 5 wandering hazards, 50
episodes each. The baseline was killed by a hazard in every episode. Zero
timeouts for Pragma confirms decisive forward progress in all episodes.}
\label{tab:gridworld}
\end{table}

The 11 killed Gridworld episodes represent hazard configurations that intersect
the direct path regardless of decision quality. The circuit breaker operated in
WARN/SLOW range (RPN 180--200) throughout --- the system correctly self-constrained
without reaching STOP.

% ── 4. Episodic Memory ────────────────────────────────────────────────────────
\section{Episodic Memory}
\label{sec:memory}

\subsection{Implementation}

Stage 7 produces Beta posterior parameters $(a, b)$ for each tracked hazard rate.
The episodic memory module persists this state to \texttt{artifacts/<benchmark>/memory.json}
at session end and loads it at session start, seeding Beta trackers with the
loaded $(a, b)$ rather than the uniform prior $\mathrm{Beta}(1, 1)$. An optional
decay factor $\alpha \in (0, 1]$ shrinks accumulated pseudo-counts toward the
uniform prior:
\begin{equation}
  a_{\mathrm{new}} = 1 + (a_{\mathrm{loaded}} - 1) \cdot \alpha
  \label{eq:decay}
\end{equation}

We additionally implemented $p_{\mathrm{death}}$ blending: at each step, the
Monte Carlo estimate is blended with the tracker's current mean:
\begin{equation}
  p_{\mathrm{death,adj}} = 0.7 \cdot p_{\mathrm{death}}^{\mathrm{MC}}
                         + 0.3 \cdot \mu_{\mathrm{tracker}}
  \label{eq:blend}
\end{equation}
The adjusted value propagates through FMEA, circuit breaker, utility, and the
subsequent belief update.

\subsection{Experiment and Results}

A two-pass comparison was run on all three benchmarks (50 episodes per pass,
seeds 0--49). Pass~1: no memory file present (uniform priors, cold start).
Pass~2: memory loaded from Pass~1 (warm start). Both prior seeding alone and
the blended $p_{\mathrm{death}}$ were tested.

\textbf{Result:} zero measurable difference across all metrics in all three
benchmarks under both implementations.

\subsection{Analysis}

Three structural reasons account for the null result.

\textbf{(i) Deterministic seeds remove stochasticity between passes.}
Both passes use \texttt{seed=0..49} on identical environment instances.
Monte Carlo rollouts use the episode seed as \texttt{seed\_base}, so
$p_{\mathrm{death}}^{\mathrm{MC}}$ is identical in both passes for every action
at every step. The blended value can only differ by the prior term, which is
itself derived from the same Pass~1 run.

\textbf{(ii) In-session accumulation floods the inter-session prior.}
Within 50 episodes, Beta trackers accumulate thousands of pseudo-counts
(e.g.\ $\mathrm{Beta}(10028, 1)$ for \texttt{death\_rate} in Snake). A loaded
prior of the same magnitude is immediately overwhelmed by in-session updates
converging to the same posterior. The prior's influence vanishes within the
first 3--4 episodes of Pass~2.

\textbf{(iii) Performance ceilings.} Maze achieves 50/50 and Snake averages
22.5 on deterministic seeds regardless of prior calibration. There is no
headroom for memory-enabled lift.

Episodic memory is an adaptation mechanism, not a performance mechanism. Its
value is proportional to: (a) stochasticity between sessions, (b) a cold-start
period where prior calibration can materially change early decisions, and
(c) available headroom below performance ceilings. These conditions are absent
in fixed-seed 50-episode benchmarks but are present in real deployments with
drifting environments and short session budgets.

% ── 5. Lessons Learned ────────────────────────────────────────────────────────
\section{Lessons Learned}
\label{sec:lessons}

\paragraph{L1 --- The safety pipeline is not the performance bottleneck.}
Across all three benchmarks and all calibration iterations, the FMEA gate and
circuit breaker operated correctly without modification. Every performance failure
was traceable to the utility function or critical path estimate. The pipeline is
a stable, domain-agnostic layer. Improving it did not improve performance;
improving the utility function and signal calibration did.

\paragraph{L2 --- The utility function is the primary design variable.}
Each benchmark improvement came from fixing the goal-progress signal, not the
safety constraints. Snake improved 57$\times$ by increasing the food-distance
weight from 0.2 to 1.5. Maze improved from 6\% to 100\% by replacing Manhattan
distance with BFS path distance. The correct architectural separation is
\textit{invariant safety} (DIC stages 1--5) and \textit{domain-specific progress
signal} (Stage~6 utility). Safety should remain fixed; the utility function is
tuned per domain.

\paragraph{L3 --- Critical path calibration must respect domain solution length.}
In Maze v1.0, rollout depth (25 steps) was shorter than the minimum maze solution
length under the rollout policy. $p_{\mathrm{death}}$ was never triggered,
silently producing zero for every candidate. The correct principle: calibrate
depth to the \textit{expected time-to-failure under the rollout policy}, not to
the episode step budget.

\paragraph{L4 --- Signal accuracy without variance is uninformative.}
After the Maze v1.0 calibration fix, $p_{\mathrm{death}} \approx 1.0$ for all
candidate actions --- accurate, but providing no basis for action selection. Useful
risk estimation requires variance across the action set, not merely correctness in
aggregate. The Gridworld environment is the first in this suite where
$p_{\mathrm{death}}$ is genuinely differentiating: moving toward a hazard cluster
has measurably higher $p_{\mathrm{death}}$ than WAIT or evasion.

% ── 6. Conclusion ─────────────────────────────────────────────────────────────
\section{Conclusion}
\label{sec:conclusion}

We presented the Decision Intelligence Core --- a seven-stage pre-execution
governance layer for agentic AI systems --- and validated it empirically across
three benchmark environments against a controlled random baseline. The DIC achieved
100\% solve rate in a maze environment where a random policy solved 6\%, 78\%
survival in a dynamic hazard environment where a random policy survived 0\%, and
consistent $+$780--1081\% improvement over baseline in Snake.

The safety pipeline was stable across all conditions. Performance was governed by
utility function design and critical path calibration. Episodic memory via Bayesian
prior transfer is architecturally sound but requires stochastic, variable, or
short-session environments to produce measurable improvement over in-session
Bayesian updating alone.

The DIC enforces the distinction between \textit{intelligence that acts} and
\textit{intelligence that acts within justifiable bounds} --- making the latter a
property of the system rather than a property of any particular upstream model.

\bigskip
\noindent\textbf{Code and replication.}
All benchmark environments, agent implementations, and experiment runners are
available at \url{https://github.com/zabinskirafal/AGI-Pragma-Core}.

% ── End ───────────────────────────────────────────────────────────────────────
\end{document}
